% Shared main-text experimental protocol; results are prospective.
\section{Experiments}
\label{sec:setup}

We optimize a harness separately in two long-input environments and then
evaluate the frozen result against three fixed-harness references. Runner,
attributor, and proposer use the same frozen \textbf{DeepSeek-V4-Flash} model,
served through OpenRouter with identical decoding settings. Optimization and
evaluation are implemented. This is a work in progress, and results and
analysis will be added once the runs complete. Preliminary runs from an earlier
configuration are reported in Appendix~\ref{sec:diag} and are not part of this
study.

\paragraph{Environments.}
\label{sec:envs-and-splits}
\emph{OOLONG-Pairs} is reconstructed from the \texttt{trec\_coarse} portion of
OOLONG-synth~\citep{bertsch2025oolong}. Each task applies one of twenty
predicates to a collection of labeled user records and asks for the set of
user-ID pairs that satisfy it. A deterministic scorer compares predicted and
gold pair sets. \emph{OBLIQ-Bench Math} (analogue queries) gives a competition
problem and a candidate corpus of other problems. The model must rank, by ID,
the candidates whose solutions use the same underlying proof technique. It is
scored by NDCG@10 against the dataset's binary relevance judgments (151 queries
over a 3,508-problem corpus).

\paragraph{Blended optimization.}
Both environments are optimized on a blend of three input sizes, so a single
harness must work across them. OOLONG-Pairs uses record collections at nominal
context windows of 8K, 16K, and 32K tokens. OBLIQ uses candidate pools of 200,
500, and 1,000 problems; each pool holds a query's gold problems plus random
negatives. At every size, mining and validation tasks are disjoint and are
never used for testing. Promotion strata $\ell$ in
Equation~\ref{eq:promotion-main} are these sizes: context windows for OOLONG-Pairs
and pool sizes for OBLIQ. Table~\ref{tab:study-settings} gives both
profiles. Both start from the sparse \hzero{}; the hand-written \hzerostar{}
is used only as an evaluation reference.

\begin{table}[t]
\centering\small
\setlength{\tabcolsep}{4pt}
\begin{tabularx}{\linewidth}{@{}l>{\raggedright\arraybackslash}Xl@{}}
\toprule
 & OOLONG-Pairs & OBLIQ Math \\
\midrule
Strata $\ell$ & 8K / 16K / 32K tok. & 200 / 500 / 1000 docs \\
Mine/val./test & 8 / 8 / 4 per $\ell$ & 8 / 8 / 4 per $\ell$ \\
Long test & 40 tasks at 262K & -- \\
$|V_s|$ & 6 (1/2/3 per $\ell$) & 6 (2 per $\ell$) \\
$m$, $v$, $v_s$, $k$ & 2, 3, 1, 4 & 2, 3, 1, 4 \\
$T$, $p$ & 10, 3 & 10, 5 \\
Quality & pair-set F1 & NDCG@10 \\
Pass & exact pair set & NDCG@10 $\ge 0.9$ \\
$\tau_\ell$ & .144 / .144 / .244 & .227 / .198 / .266 \\
$\tau_{\mathrm{imp}}$ & 0 & 0.01 \\
Cost band & $[0.5, 3]\times$ & $[0, 3]\times$ \\
Early stop & on & off \\
Re-attrib. & on & off \\
Refinement & on (\texttt{wrong\_format}) & off \\
Iter./depth & 45 (+3 final.) / 3 & 30 / 3 \\
\bottomrule
\end{tabularx}
\caption{Optimization profiles. Per-size counts are tasks; $m$ and $v$ are
attempts per mining and validation task, $v_s$ per screening task, $k$ the
candidate ceiling, $T$ the round ceiling, and $p$ the patience. Margins $\tau_\ell$ come from a bootstrap over each
environment's round-one incumbent validation scores. Both profiles cap each run
at \$2 and 3,600\,s and each validation subject at \$100.}
\label{tab:study-settings}
\end{table}

\paragraph{Evaluation conditions.}
\label{sec:baselines}
After optimization stops, the final incumbent is frozen as \shrlm and compared
on held-out test tasks with three fixed references. All four conditions run on
the same model:
\begin{itemize}
\setlength{\itemsep}{1pt}
\item \hzero{}: the sparse starting harness, measuring what optimization added;
\item \hzerostar{}: an upstream-style hand-written \rlm{} prompt;
\item $\lambda$-RLM: a pinned reconstruction of the hand-designed recursive
method of~\citet{roy2026ycombinator};
\item \shrlm: the frozen optimized harness.
\end{itemize}
Each condition makes three attempts per test task. Test tasks include the four
held-out tasks per optimization size in each environment. For OOLONG-Pairs they
also include 40 tasks at a nominal 262K window. That is $8\times$ the longest
optimization context, and no long-window task is mined or validated. OBLIQ test
queries come from a frozen 40-query split already used for earlier
full-corpus comparisons, and they are disjoint from all mining and validation
queries.

\paragraph{Metrics.}
We report mean primary quality (F1 or NDCG@10) per test size and macro-averaged
across sizes, together with pass rate under each environment's pass threshold.
We also report cost, runtime, sub-call counts, and completed, failed, and
missing attempts. Repeated attempts on one task, and OOLONG-Pairs predicates
that share a record collection, are not independent samples, so uncertainty
will be estimated over tasks. A promotion during optimization is a selection
decision, not a significance test. Cross-environment transfer and
weight-training comparisons remain future work.
